% TEMPLATE-MILC-v3.tex - plantilla maestra UNICA de las guias MILC v3.
%
% La usan las tres familias de guias, cada una con su driver:
%   · Grados 6-11  -> scripts/build-guias-g11.py
%   · Bebras       -> scripts/build-guias-bebras.py
%   · Semillero    -> scripts/build-guias-semillero.py
%
% Antes habia tres copias casi identicas (~400 lineas iguales, 6 distintas):
% cada mejora habia que aplicarla tres veces, y en la practica no se hacia
% (el motor de recursos vivia solo en dos de ellas). Lo que varia entre
% familias esta parametrizado en 6 placeholders que cada driver llena
% (se nombran aqui SIN su sintaxis real: el builder reemplaza por texto plano,
% tambien dentro de los comentarios, y un valor multilinea rompe el preambulo):
%
%   HEADER_IZQ        encabezado izquierdo de pagina
%   HEADER_DER        encabezado derecho de pagina
%   PORTADA_ETIQUETA  rotulo sobre el numero grande (GRADO/MOMENTO/MODULO)
%   PORTADA_SERIE     linea de serie de la portada
%   PORTADA_ID        identificador de la guia en la portada
%   PORTADA_CIRCULO   el circulito de grado (°), o vacio si no aplica
%
% El resto de claves las documenta content/guias/_SCHEMA.md. El script valida
% que no queden placeholders sin reemplazar antes de escribir el archivo final.

\documentclass[11pt,letterpaper]{article}
\usepackage[margin=1.75cm]{geometry}
\usepackage[table]{xcolor}
\usepackage{fontspec}
\usepackage[spanish]{babel}
\usepackage{tikz}
% Librerías TikZ para los diagramas de las guías (bloque `recursos:`):
%  · babel  → babel en español vuelve activos caracteres como «>», y TikZ los usa
%             en sus opciones (>=stealth, ->). Sin esta librería, cualquier
%             diagrama con flechas revienta con «\language@active@arg> has an
%             extra }». Debe ir ANTES de que se dibuje nada.
%  · positioning        → sintaxis «below left=2pt and 3pt».
%  · decorations.pathreplacing → llaves ({brace}) para señalar rangos.
\usetikzlibrary{babel,positioning,decorations.pathreplacing}
\usepackage{graphicx}
% Circuitos: el trabajo de electrónica se hace hoy con micro:bit y sus sensores
% integrados (MakeCode, sin cableado), así que no se necesitan esquemáticos.
% Si algún día entran montajes con protoboard, basta descomentar esta línea:
% los diagramas .tex/.tikz declarados en `recursos:` ya se incrustan solos.
% \usepackage{circuitikz}
\usepackage[most]{tcolorbox}
\usepackage{tabularx}
\usepackage{array}
\usepackage{fancyhdr}
\usepackage{titlesec}
\usepackage{ragged2e}
\usepackage{enumitem}
\usepackage{microtype}

% Helvetica Neue no trae la flecha U+2192, asi que cada «→» escrito literalmente
% en el YAML se compilaba a NADA, en silencio (462 flechas en 61 guias: rutas del
% tipo «paleta → tipografia → lockup» salian rotas). Mapeamos el caracter a su
% version matematica, que si tiene glifo. Asi el autor puede seguir escribiendo
% «→» comodamente en el YAML.
\usepackage{newunicodechar}
\newunicodechar{→}{\ensuremath{\rightarrow}}

\setmainfont{Helvetica Neue}
\setsansfont{Helvetica Neue}
\setlength{\parindent}{0pt}
\setlength{\parskip}{6pt}
\hyphenpenalty=200
\exhyphenpenalty=200
\pretolerance=400
\tolerance=2000
\setlength{\emergencystretch}{1em}
\renewcommand{\arraystretch}{1.25}
\setlist{leftmargin=5mm,itemsep=1mm,topsep=1mm}
\newcolumntype{Y}{>{\RaggedRight\arraybackslash}X}

\definecolor{milcMagenta}{HTML}{E6007E}
\definecolor{milcVino}{HTML}{5A0038}
\definecolor{milcTurquesa}{HTML}{0093A5}
\definecolor{milcVerde}{HTML}{58A923}
\definecolor{milcMostaza}{HTML}{E5B400}
\definecolor{milcOcre}{HTML}{B86600}
\definecolor{milcNegro}{HTML}{241816}
\definecolor{milcGris}{HTML}{F7F7F4}
\definecolor{milcLinea}{HTML}{E8E2DD}
\definecolor{milcGradePrimary}{HTML}{FF2D87}
\definecolor{milcGradeDark}{HTML}{9F1B56}
\definecolor{milcGradeSoft}{HTML}{FFE0EE}
\definecolor{milcGradeAccent}{HTML}{CFE7FF}
\definecolor{milcGradeText}{HTML}{FFFFFF}
\color{milcNegro}

\pagestyle{fancy}
\fancyhf{}
\lhead{\small Tecnología e Informática · Grado 10}
\rhead{\small Guía 16 · MILC v3}
\cfoot{\small\thepage}
\renewcommand{\headrulewidth}{0pt}

\newtcolorbox{titlebox}[1]{
  enhanced,colback=#1,colframe=#1,arc=7mm,boxrule=0pt,
  left=7mm,right=7mm,top=3mm,bottom=3mm,
  before skip=8pt,after skip=8pt,
  fontupper=\sffamily\bfseries\Large\color{white}
}

\newtcolorbox{softbox}[3]{
  enhanced,colback=#2,colframe=#1,borderline west={4pt}{0pt}{#1},
  arc=2mm,boxrule=.4pt,left=4mm,right=4mm,top=3mm,bottom=3mm,
  before skip=6pt,after skip=8pt,title={#3},coltitle=white,
  colbacktitle=#1,fonttitle=\sffamily\bfseries,fontupper=\RaggedRight,
  attach boxed title to top left={xshift=3mm,yshift=-2mm},
  boxed title style={arc=2mm, boxrule=0pt}
}

\newtcolorbox{drawbox}[2]{
  enhanced,colback=white,colframe=#1,arc=4mm,boxrule=1pt,
  left=5mm,right=5mm,top=4mm,bottom=4mm,before skip=8pt,after skip=8pt,
  title={#2},coltitle=white,colbacktitle=#1,fonttitle=\sffamily\bfseries,
  fontupper=\RaggedRight,
  attach boxed title to top left={xshift=4mm,yshift=-2mm},
  boxed title style={arc=3mm, boxrule=0pt}
}

\newtcolorbox{infoband}[1]{
  enhanced,colback=#1!8,colframe=#1!50,arc=2mm,boxrule=.5pt,
  left=4mm,right=4mm,top=2mm,bottom=2mm,before skip=4pt,after skip=4pt,
  fontupper=\small\RaggedRight
}

% Puente narrativo entre secciones (contrato editorial Conéctate).
% Da continuidad: "Acabas de X. Ahora vas a Y porque Z."
% Visualmente: banda izquierda en color del grado, texto en itálica pequeña.
\newtcolorbox{puentebox}{
  enhanced,colback=milcGradeSoft,colframe=milcGradeDark,
  borderline west={3pt}{0pt}{milcGradeDark},
  arc=0mm,boxrule=0pt,left=5mm,right=4mm,top=2.5mm,bottom=2.5mm,
  before skip=4pt,after skip=8pt,
  fontupper=\itshape\small\color{milcNegro}\RaggedRight
}
% La flecha va en modo matematico a proposito: Helvetica Neue NO tiene el glifo
% U+2192, asi que \textrightarrow se compilaba a NADA («Missing character: There
% is no → in font Helvetica Neue») y el puente salia sin flecha. En modo
% matematico la flecha viene de la fuente matematica y si se dibuja.
\newcommand{\puente}[1]{%
  \begin{puentebox}%
    \textcolor{milcGradeDark}{$\rightarrow$}\hspace{2mm}#1%
  \end{puentebox}%
}

\newcommand{\linead}[1]{\noindent\color{milcLinea}\rule{#1}{0.5pt}\color{milcNegro}}
\newcommand{\respuesta}[1]{\par\vspace{2mm}\foreach \n in {1,...,#1}{\noindent\color{milcLinea}\rule{\linewidth}{0.5pt}\par\vspace{5mm}}\color{milcNegro}}
\newcommand{\checkbox}{\textcolor{milcGradeDark}{\fbox{\phantom{x}}}\hspace{5pt}}

\newcommand{\phasecard}[4]{
\begin{tcolorbox}[enhanced,colback=#3,colframe=#2,arc=3mm,boxrule=.5pt,height=3.05cm,valign=center,halign=center,left=1.5mm,right=1.5mm,top=2mm,bottom=2mm]
{\sffamily\bfseries\fontsize{22}{24}\selectfont\textcolor{#2}{#1}}\par
{\sffamily\bfseries\small #4}
\end{tcolorbox}
}

% ── Recursos visuales de la guía ──────────────────────────────────────
% Declarados en el bloque `recursos:` del YAML y emitidos por el builder.
% \guiaFigura   → imagen rasterizada o PDF (foto, carta celeste, montaje).
% \guiaDiagrama → fuente TikZ/circuitikz (.tex) incrustada como vector:
%                 es la vía para circuitos y esquemáticos editables.
% #1 = ruta relativa al .tex · #2 = pie (puede ir vacío)
\newcommand{\guiaFigura}[2]{%
\begin{center}
\includegraphics[width=0.88\linewidth,height=0.38\textheight,keepaspectratio]{#1}\\[1.5mm]
{\footnotesize\itshape\textcolor{milcNegro!75}{#2}}
\end{center}
}

\newcommand{\guiaDiagrama}[2]{%
\begin{center}
\input{#1}\\[1.5mm]
{\footnotesize\itshape\textcolor{milcNegro!75}{#2}}
\end{center}
}

\begin{document}
\thispagestyle{empty}
\begin{tikzpicture}[remember picture,overlay]
  \fill[white] (current page.south west) rectangle (current page.north east);
  \path[fill=milcGradePrimary,rounded corners=16mm]
    ([xshift=.55cm,yshift=-.55cm]current page.north west) rectangle
    ([xshift=-.55cm,yshift=3.65cm]current page.south east);

  \node[anchor=north west,text=milcGradeText] at ([xshift=2.0cm,yshift=-1.85cm]current page.north west)
    {\sffamily\bfseries\fontsize{14}{16}\selectfont GRADO};
  \node[anchor=north west,text=milcGradeText,opacity=.82] at ([xshift=2.0cm,yshift=-2.75cm]current page.north west)
    {\sffamily\bfseries\fontsize{9}{11}\selectfont SERIE GUÍAS · TECNOLOGÍA E INFORMÁTICA};

  \begin{scope}[shift={([xshift=-3.05cm,yshift=-2.55cm]current page.north east)}]
    \fill[white,opacity=.80] (-.62,-.52) rectangle (.62,.52);
    \draw[milcGradeText,line width=1pt] (-.62,-.52) rectangle (.62,.52);
    \fill[milcGradeDark] (-.42,-.38) rectangle (-.22,.06);
    \fill[milcGradeAccent] (-.10,-.38) rectangle (.10,.24);
    \fill[milcGradePrimary!60!black] (.22,-.38) rectangle (.42,.38);
  \end{scope}

  \node[anchor=west,text=milcGradeText] at ([xshift=1.9cm,yshift=-12.85cm]current page.north west)
    {\sffamily\bfseries\fontsize{124}{124}\selectfont 10};
  % El circulito de grado (°) solo aplica a las guias de grado («11°»); en
  % Bebras y Semillero el numero grande es un momento/modulo, no un grado.
  \draw[milcGradeText,line width=1.7pt]
    ([xshift=4.52cm,yshift=-11.68cm]current page.north west) circle (.13cm);

  \node[anchor=west,text=milcGradeText,text width=8.7cm,align=left] at ([xshift=10.35cm,yshift=-11.6cm]current page.north west)
    {
     {\sffamily\bfseries\fontsize{19}{22}\selectfont Análisis cualitativo\\con IA ---\\qué hace y qué no}\\[4mm]
     {\sffamily\bfseries\fontsize{23}{26}\selectfont Guía 16-10-TIC}\\[2mm]
     {\sffamily\fontsize{13}{16}\selectfont Período 2 · Informes técnicos y comunicación profesional con IA}\\[5mm]
     {\sffamily\bfseries\fontsize{11}{13}\selectfont Institución Educativa Sor María Juliana}\\[1mm]
     {\sffamily\fontsize{10}{12}\selectfont Autor: Dr. Álvaro Cárdenas Orozco}
    };

  \path[fill=milcGradeSoft,rounded corners=7mm]
    ([xshift=1.35cm,yshift=.85cm]current page.south west) rectangle
    ([xshift=-1.35cm,yshift=4.25cm]current page.south east);
  \draw[milcGradeDark,line width=.65pt,rounded corners=7mm]
    ([xshift=1.35cm,yshift=.85cm]current page.south west) rectangle
    ([xshift=-1.35cm,yshift=4.25cm]current page.south east);
  \node[anchor=center,text=milcNegro,text width=17.0cm,align=center] at ([yshift=2.55cm]current page.south)
    {\sffamily\fontsize{10}{12}\selectfont Metodología MILC · Apertura ancestral · Pensamiento computacional · Triángulo Dussel-Estoicismo-Floridi\\
    \textcolor{milcGradeDark}{\bfseries Producto:} Análisis cualitativo con asistencia de IA de 15-20 respuestas abiertas reales de tu encuesta de la sesión 5 + clasificación temática automática propuesta por la IA (3-5 categorías emergentes) + tabla de auditoría con 3 columnas (a) lo que la IA clasificó correctamente, (b) lo que erró o malinterpretó, (c) cómo lo corregiste con criterio humano + reflexión final de 5 líneas sobre los límites del análisis automatizado con respuestas en español coloquial};
\end{tikzpicture}
\mbox{}
\newpage

\section*{Apertura: Análisis cualitativo con IA --- qué hace y qué no}

\begin{softbox}{milcGradeDark}{milcGradeSoft}{Saber ancestral}
En la cocina de los restaurantes pequeños del centro de Cartago, sostuvo durante décadas una práctica que cualquier comensal observador reconoce: \textbf{la cocinera muestra siempre las medidas antes de cocinar}. Cuando un cliente pide un plato especial o complejo, la cocinera no entra a la cocina directamente: saca la libra de pollo, los gramos de arroz, la taza de caldo, los pone sobre la mesa, y dice: \emph{``vea, esto es lo que va''}. El cliente mira, aprueba o sugiere ajuste (\emph{``menos sal''}, \emph{``más papa''}), y solo entonces la cocinera entra a la estufa. Ese gesto cumple 2 funciones simultáneas: (1) \textbf{Transparencia}: el cliente ve lo que va a comer antes de que esté cocinado. (2) \textbf{Verificación humana}: el cliente puede corregir antes del punto sin retorno. La sabiduría es ancestral: \textbf{quien cocina y quien come deciden juntos, no la cocinera sola}. Esa práctica del oficio se traduce directamente al análisis de datos con IA: la IA es buena agregando y proponiendo clasificaciones, pero puede errar en matices culturales o coloquiales. Tu papel como editor es \textbf{ver lo que propone la IA antes de servirlo como conclusión}.
\end{softbox}

\begin{softbox}{milcTurquesa}{milcGris}{Saber contemporáneo}
El \textbf{análisis cualitativo} es el proceso de extraer significado de respuestas abiertas (texto, no números). Tradicionalmente requería leer cada respuesta, agrupar manualmente por temas, codificar con etiquetas. Con IA generativa (ChatGPT, Gemini, Claude), el proceso se acelera: pegas las respuestas y le pides clasificación temática automática. La IA propone categorías emergentes (\emph{salud mental, presión académica, problemas de transporte, etc.}) y asigna cada respuesta a una categoría. La regla profesional famosa: \textbf{la IA propone, el humano dispone}. La IA es buena en: (a) detectar patrones obvios. (b) agrupar respuestas similares. (c) hacer resúmenes. La IA es mala en: (a) detectar sarcasmo (\emph{``perfecto, otro recreo de 5 minutos''} puede tomarse literal). (b) entender modismos locales (\emph{``estoy planchao''} = sin dinero en costa, no en interior). (c) captar contexto cultural específico. (d) detectar respuestas que parecen claras pero esconden ambigüedad. El editor humano debe \textbf{auditar la clasificación de la IA siempre}, especialmente con respuestas en español coloquial colombiano.
\end{softbox}

\begin{softbox}{milcMostaza}{milcGris}{Pregunta-puente}
\textit{¿Qué sabía la cocinera al mostrar las medidas antes de cocinar, que el novato olvida cuando acepta la clasificación de la IA sin verificar? ¿Y por qué la IA puede equivocarse interpretando un \emph{``más o menos''} colombiano que es ambiguo a propósito?}
\end{softbox}

\begin{softbox}{milcVerde}{milcGris}{Saber y saber hacer}
Al terminar podrás: (1) \textbf{identificar} qué tipos de respuesta son fáciles para la IA y cuáles le dan problema; (2) \textbf{analizar} 15-20 respuestas abiertas con asistencia de IA para clasificación temática; (3) \textbf{evaluar} la clasificación de la IA respuesta por respuesta, detectando aciertos y errores; (4) \textbf{explicar} con honestidad los límites del análisis automatizado y cómo el criterio humano sigue siendo necesario.
\end{softbox}



\small
\begin{tabularx}{\textwidth}{>{\bfseries}p{3.0cm}Y}
\rowcolor{milcGradePrimary!12}Autor & Dr. Álvaro Cárdenas Orozco\\
DBA & Comunicación profesional con asistencia de IA y herramientas ofimáticas gratuitas (MEN, T\&I)\\
Referentes & Markdown como estándar abierto · Google Workspace · Estoicismo (claridad)\\
Producto final & Análisis cualitativo con asistencia de IA de 15-20 respuestas abiertas reales de tu encuesta de la sesión 5 + clasificación temática automática propuesta por la IA (3-5 categorías emergentes) + tabla de auditoría con 3 columnas (a) lo que la IA clasificó correctamente, (b) lo que erró o malinterpretó, (c) cómo lo corregiste con criterio humano + reflexión final de 5 líneas sobre los límites del análisis automatizado con respuestas en español coloquial\\
\end{tabularx}
\normalsize

\puente{Antes de empezar, mira el viaje completo. En la sesión 5 diseñaste encuesta con filtro de IA. Hoy aprendes a \emph{analizar los resultados} con IA pero con criterio humano. Las próximas sesiones (informe comercial, correo profesional, mini-proyecto, sustentación) usarán los hallazgos del análisis de hoy.}

\begin{titlebox}{milcGradeDark}
Ruta MILC: así vamos a aprender
\end{titlebox}
\begin{tabularx}{\textwidth}{>{\centering\arraybackslash}X>{\centering\arraybackslash}X>{\centering\arraybackslash}X>{\centering\arraybackslash}X}
\phasecard{1}{milcVerde}{milcGris}{Escucha\\[-1mm]\small Observo, escucho y pregunto.} &
\phasecard{2}{milcTurquesa}{milcGris}{Sistematización\\[-1mm]\small Pensamiento computacional.} &
\phasecard{3}{milcMagenta}{milcGris}{Praxis\\[-1mm]\small Construyo y aplico.} &
\phasecard{4}{milcVino}{milcGris}{Evaluación\\[-1mm]\small Reflexión liberadora.}
\end{tabularx}


\puente{Antes de teorizar, vas a hacer una \emph{prueba simple}: toma 5 respuestas abiertas de tu encuesta y léelas tú mismo a mano. Sin IA todavía. Anota qué temas detectas. Después harás el mismo análisis con IA y compararás resultados.}

\begin{titlebox}{milcVerde}
Fase 1 · Escucha
\end{titlebox}

\begin{softbox}{milcVerde}{milcGris}{Escena para escuchar}
\textbf{Actividad 1 · IDENTIFICA --- Lectura humana de 5 respuestas} (15 min · individual). Toma 5 respuestas abiertas de tu encuesta (sesión 5). Léelas con calma a mano (sin IA). Anota: (1) ¿qué tema dominante detectas en cada una? (2) ¿alguna tiene ironía o sarcasmo? (3) ¿alguna usa modismos colombianos que requieren contexto? (4) ¿qué clasificación temática propondrías para las 5? Esa lectura humana será tu referencia para comparar después con la IA.
\end{softbox}

\checkbox Leí 5 respuestas a mano sin IA.\\
\checkbox Detecté tema, ironía, modismos.\\
\checkbox Propuse mi propia clasificación.

\begin{infoband}{milcVerde}


\textbf{Lo que va al cuaderno (Actividad 1 · IDENTIFICA).} Título: «Actividad 1 --- Mi lectura humana». \textbf{Formato:} 5 fichas (1 por respuesta) + clasificación propuesta a mano. \textbf{Sabes que terminaste cuando} tienes referencia humana para comparar con IA.
\end{infoband}

\puente{Ya tienes tu lectura humana. Ahora vas a entender la \textbf{anatomía del análisis cualitativo con IA}: cómo pedirle clasificación, cómo verificar, los 5 errores típicos de aceptar IA sin auditoría.}

\begin{titlebox}{milcTurquesa}
Fase 2 · Sistematización con pensamiento computacional
\end{titlebox}

\begin{softbox}{milcTurquesa}{milcGris}{Cuatro pilares aplicados al tema}
Tu lectura humana es la referencia. Ahora necesitas la \textbf{anatomía profesional} del análisis cualitativo con IA y sus límites.
\end{softbox}

\small
\begin{tabularx}{\textwidth}{>{\bfseries\RaggedRight\arraybackslash}p{3.6cm}Y}
\rowcolor{milcTurquesa!90}\color{white}Pilar & \color{white}\bfseries Aplicación al tema\\
1. Descomposición & El \textbf{análisis cualitativo con IA} sigue un proceso estándar: (1) \textbf{Preparar las respuestas}: copiar todas las respuestas abiertas en una sola lista numerada, anonimizadas si es necesario. (2) \textbf{Prompt de clasificación}: \emph{``Actúa como analista cualitativo. Aquí están 15 respuestas a la pregunta [X] de una encuesta a estudiantes de grado 10 en Colombia. Analízalas y propone (a) 3-5 categorías temáticas emergentes, (b) asigna cada respuesta a una categoría, (c) identifica patrones, (d) señala respuestas ambiguas que requieren revisión humana''}. (3) \textbf{Aplicar prompt}: en ChatGPT, Gemini o Claude. (4) \textbf{Auditar la clasificación}: comparar respuesta por respuesta tu lectura humana con la propuesta de la IA. (5) \textbf{Corregir}: ajustar categorías o reclasificar respuestas que la IA malinterpretó.\\
2. Reconocimiento de patrones & Los \textbf{5 errores típicos del análisis con IA} sin auditoría: (a) \textbf{IA toma sarcasmo como literal}: \emph{``buenísimo el recreo de 5 minutos''} clasificado como satisfacción. (b) \textbf{IA no entiende modismos locales}: \emph{``está chévere''} clasificado neutro cuando es positivo. (c) \textbf{IA inventa patrones que no están}: alucinación de tendencias. (d) \textbf{IA genera categorías muy amplias o muy estrechas}: a veces 2 categorías para 20 respuestas (demasiado amplio), a veces 12 (demasiado fragmentado). (e) \textbf{Aceptar el resumen sin verificar las respuestas individuales}: el resumen puede ser bonito pero no reflejar las respuestas reales.\\
3. Abstracción & La \textbf{tabla de auditoría} con 3 columnas se construye así: (a) \textbf{Respuesta} (texto resumido, número de orden). (b) \textbf{Clasificación de la IA}. (c) \textbf{Tu juicio}: \emph{correcto}, \emph{incorrecto}, \emph{ambiguo}. (d) \textbf{Si incorrecto o ambiguo}: tu clasificación corregida + razón breve. Esta auditoría demuestra que la IA fue herramienta, no fuente única.\\
4. Algoritmo conceptual & Algoritmo: (1) preparar 15-20 respuestas anonimizadas. (2) escribir prompt de clasificación. (3) pasar a IA. (4) revisar la clasificación respuesta por respuesta. (5) llenar tabla de auditoría. (6) corregir errores. (7) escribir reflexión sobre límites.\\
\end{tabularx}
\normalsize

\begin{drawbox}{milcTurquesa}{Anatomía del análisis cualitativo con IA}
\textbf{Preparar:} respuestas anonimizadas. \\[1mm] \textbf{Prompt:} rol analista + clasificación pedida. \\[1mm] \textbf{Auditar:} respuesta por respuesta. \\[1mm] \textbf{Corregir:} reclasificar errores. \\[1mm] \textbf{Reflexión:} límites del análisis IA.
\end{drawbox}

\begin{softbox}{milcMostaza}{milcGris}{Errores comunes y su corrección}
(1) IA toma sarcasmo como literal. (2) IA no entiende modismos locales. (3) IA alucina patrones que no están. (4) IA propone categorías muy amplias o estrechas. (5) Aceptar resumen sin verificar respuestas. (6) Sin tabla de auditoría: la IA queda como fuente única. (7) No declarar uso de IA en el informe final.
\end{softbox}

\begin{infoband}{milcTurquesa}


\textbf{Lo que va al cuaderno (Actividad 2 · ANALIZA).} Título: «Actividad 2 --- Anatomía análisis con IA». \textbf{Formato:} (a) los 5 pasos del proceso; (b) los 7 errores con marca; (c) plantilla de la tabla de auditoría. \textbf{Sabes que terminaste cuando} podrías auditar análisis ajeno con criterio.
\end{infoband}

\puente{Tienes el método. Toca aplicarlo: análisis con IA de 15-20 respuestas, comparación con tu lectura humana, tabla de auditoría con 3 columnas, reflexión sobre límites.}

\begin{titlebox}{milcMagenta}
Fase 3 · Praxis con pensamiento computacional
\end{titlebox}

\begin{softbox}{milcMagenta}{milcGris}{Construyendo el producto con los cuatro pilares}
\textbf{Actividad 3 · EVALÚA + EXPLICA --- Análisis IA + tabla auditoría + reflexión} (40 min · individual con dispositivo). Hora de aplicar. Análisis con IA de 15-20 respuestas, comparación con tu lectura humana de la Actividad 1, tabla de auditoría completa, reflexión sobre límites.
\end{softbox}

\small
\begin{tabularx}{\textwidth}{>{\bfseries\RaggedRight\arraybackslash}p{3.6cm}Y}
\rowcolor{milcMagenta!90}\color{white}Pilar & \color{white}\bfseries Aplicación al producto\\
Descomposición & Tu análisis se construye en 5 pasos: (1) tomar 15-20 respuestas abiertas de tu encuesta (incluir las 5 que ya leíste a mano). (2) escribir prompt de clasificación con rol + contexto + tarea específica. (3) ejecutar en ChatGPT, Gemini o Claude. (4) construir tabla de auditoría comparando lectura humana vs IA respuesta por respuesta. (5) escribir reflexión final.\\
Patrones & Sigue el patrón \textbf{``confiar pero verificar''}: la IA es asistente útil pero no autoridad. Cada categoría que propone y cada asignación que hace debe pasar tu auditoría. La phronesis del oficio consiste en \textbf{usar la IA para acelerar el análisis, no para reemplazarlo}.\\
Abstracción & La reflexión final de 5 líneas responde: (a) ¿cuántas respuestas la IA clasificó correctamente? (b) ¿en qué tipos de respuestas erró más? (c) ¿qué modismos colombianos confundieron a la IA? (d) ¿qué cambiarías de tu prompt para mejorar? (e) ¿cuándo conviene IA y cuándo solo lectura humana?\\
Algoritmo de armado & Algoritmo: (1) preparar respuestas. (2) prompt. (3) IA. (4) tabla auditoría. (5) corregir. (6) reflexión. (7) declarar uso de IA en cualquier informe que use estos hallazgos.\\
\end{tabularx}
\normalsize

\begin{drawbox}{milcMagenta}{Checklist antes de entregar}
\checkbox 15-20 respuestas abiertas analizadas. \\[1mm] \checkbox Prompt de clasificación bien estructurado. \\[1mm] \checkbox Clasificación IA generada. \\[1mm] \checkbox Tabla de auditoría completa con 3 columnas. \\[1mm] \checkbox Errores de la IA identificados y corregidos. \\[1mm] \checkbox Reflexión final de 5 líneas. \\[1mm] \checkbox Uso de IA declarado para futuros informes.
\end{drawbox}

\begin{softbox}{milcMagenta}{milcGris}{Plantilla de trabajo}
\textbf{Respuestas.} \textit{15-20 anonimizadas.} \\[1mm] \textbf{Prompt.} \textit{Rol + contexto + clasificación pedida.} \\[1mm] \textbf{Clasificación IA.} \textit{Categorías + asignaciones.} \\[1mm] \textbf{Tabla auditoría.} \textit{Respuesta + IA + Juicio + Corrección.} \\[1mm] \textbf{Reflexión.} \textit{Aciertos, errores, modismos, ajustes, cuándo IA y cuándo no.}
\end{softbox}

\begin{infoband}{milcMagenta}


\textbf{Lo que va al cuaderno (Actividad 3 · EVALÚA + EXPLICA).} Título: «Actividad 3 --- Análisis con criterio humano». \textbf{Formato:} (a) prompt usado; (b) tabla auditoría; (c) reflexión. \textbf{Sabes que terminaste cuando} podrías defender cada clasificación final ante el grupo.
\end{infoband}

\puente{Lo que sigue resume \emph{qué} entregas hoy: análisis IA + tabla auditoría + reflexión. El criterio principal: que la conclusión que escribas resista verificación humana paso a paso.}

\begin{titlebox}{milcMostaza}
Producto de la sesión
\end{titlebox}
\textbf{Análisis IA + tabla de auditoría + reflexión sobre límites}.
\begin{softbox}{milcMostaza}{milcGris}{Criterios de calidad}
(1) 15-20 respuestas analizadas. (2) Tabla de auditoría con 3 columnas. (3) Errores identificados y corregidos. (4) Categorías finales defendibles. (5) Reflexión honesta sobre límites. (6) Uso de IA declarado.
\end{softbox}

\puente{Antes de cerrar, mira el análisis desde las cinco dimensiones humanas. Verificar la IA respuesta por respuesta es respeto por los encuestados; aceptar la clasificación ciegamente es negligencia profesional.}

\begin{titlebox}{milcVino}
Fase 4 · Evaluación liberadora — cinco dimensiones
\end{titlebox}

\begin{softbox}{milcVino}{milcGris}{Sentido de la evaluación}
Evaluar no es solo revisar una nota. Es reconocer qué aprendí, qué cambió en mí, qué emociones manejé, cómo me relaciono con la ciudad y la comunidad, y qué le dejaré a quien viene detrás.
\end{softbox}

\textbf{Mi reflexión liberadora en cinco dimensiones.} Completa cada frase iniciada:

\small
\begin{tabularx}{\textwidth}{>{\bfseries\RaggedRight\arraybackslash}p{3.4cm}Y>{\RaggedRight\arraybackslash}p{4.6cm}}
\rowcolor{milcVino}\color{white}Dimensión & \color{white}\bfseries Mi respuesta (frame starter) & \color{white}\bfseries Algo que descubrí\\
1. Personal & Lo que más cambió en mí fue \linead{6.5cm} & \linead{4.2cm}\\
2. Emocional & La emoción más fuerte fue \linead{2.5cm} y la regulé \linead{3cm} & \linead{4.2cm}\\
3. Ciudadana & En democracia esto importa porque \linead{6cm} & \linead{4.2cm}\\
4. Local (Cartago) & En mi barrio sería útil para \linead{6.5cm} & \linead{4.2cm}\\
5. Intergeneracional & A mi abuela/abuelo le diría \linead{6.5cm} & \linead{4.2cm}\\
\end{tabularx}
\normalsize

\begin{infoband}{milcVino}
\textbf{Para la próxima sustentación.} Vuelve a esta tabla en seis meses. Verás que las respuestas cambian — no porque lo aprendido cambie, sino porque tú cambias. La evaluación liberadora es una conversación contigo a través del tiempo.
\end{infoband}


\puente{Y para terminar, tres voces sobre la verificación. Dussel sobre las clasificaciones que invisibilizan voces, Marco Aurelio sobre la disciplina de revisar, Floridi sobre la auditoría humana como ética del análisis con IA.}

\begin{titlebox}{milcGradeDark}
Triángulo de pensamiento · cierre filosófico
\end{titlebox}

\begin{softbox}{milcVino}{milcGris}{Dussel · lente del nosotros}
\textit{``Las clasificaciones automáticas pueden invisibilizar voces; la auditoría humana las rescata.''}

Cuando la IA clasifica una respuesta sarcástica como literal o un modismo coloquial como neutro, está silenciando lo que el encuestado realmente quiso decir. La filosofía de la liberación pide auditar para \textbf{no permitir que el algoritmo invisibilice voces}. Cada respuesta que la IA malinterpreta es una persona que no fue escuchada bien. La phronesis del oficio consiste en \textbf{rescatar las voces que la IA pierde}.

\textbf{¿Mi auditoría rescató las voces que la IA malinterpretó, o acepté la clasificación sin verificar?}
\end{softbox}
\linead{\linewidth}\\[1mm]
\linead{\linewidth}

\begin{softbox}{milcMostaza}{milcGris}{Estoicismo · Marco Aurelio · cuidado interior}
\textit{``Revisar respuesta por respuesta es disciplina; aceptar el resumen sin verificar es vanidad.''}

Es tentador aceptar el resumen bonito que produce la IA y pasar al siguiente paso. La sabiduría estoica recomienda lo opuesto: \emph{revisa cada respuesta antes de aceptar el resumen}. Esa disciplina parece lenta pero es la única que produce análisis confiables. La IA es buena ayudante; nunca es autoridad final.

\textbf{¿Estoy revisando con disciplina cada clasificación, o quiero saltar al resumen porque parece más rápido?}
\end{softbox}
\linead{\linewidth}\\[1mm]
\linead{\linewidth}

\begin{softbox}{milcTurquesa}{milcGris}{Floridi · lente de la infoesfera}
\textit{``La auditoría humana del análisis con IA es la nueva ética profesional del oficio analítico contemporáneo.''}

En la era de la IA generativa, miles de análisis se publican sin auditoría humana. La ética profesional pide lo opuesto: cuando usas IA para analizar, audita cada paso. Tu tabla de auditoría es ejemplo de esa práctica. Aunque sea proyecto escolar, la disciplina es la misma que rige toda análisis profesional serio con IA.

\textbf{¿Mi análisis tiene auditoría humana real, o solo pego lo que la IA me dijo?}
\end{softbox}
\linead{\linewidth}\\[1mm]
\linead{\linewidth}

\begin{titlebox}{milcVino}
Compromiso final
\end{titlebox}

\small
\begin{tabularx}{\textwidth}{>{\RaggedRight\arraybackslash}p{3.0cm}YYY}
\rowcolor{milcVino}\color{white}\bfseries Criterio & \color{white}\bfseries Lo logré & \color{white}\bfseries En proceso & \color{white}\bfseries Necesito apoyo\\
Comprensión del tema & Explico ideas centrales con mis palabras. & Reconozco ideas pero falta precisión. & Necesito repasar conceptos base.\\
Pensamiento computacional & Apliqué los 4 pilares al diseñar mi producto. & Apliqué algunos pilares. & Necesito releer la fase 2.\\
Aplicación y producto & Mi producto cumple los criterios y se puede revisar. & Producto funcional con ajustes pendientes. & Aún en construcción.\\
Reflexión 5 dimensiones & Respondí cada dimensión con honestidad. & Respondí algunas con detalle. & Mis respuestas son muy generales.\\
Triángulo de pensamiento & Conecté las 3 voces con mi trabajo real. & Conecté al menos una. & No relaciono las voces con mi proceso.\\
\end{tabularx}
\normalsize

\textbf{Hoy entendí que:}\\[1mm]
\linead{\linewidth}\\[1mm]
\linead{\linewidth}

\textbf{Mi compromiso para la próxima sesión es:}\\[1mm]
\linead{\linewidth}\\[1mm]
\linead{\linewidth}

\begin{softbox}{milcMostaza}{milcGris}{Cita de despedida · Marco Aurelio}
\textit{``El obstáculo en el camino se vuelve el camino. Lo que estorba al camino se vuelve camino.''}\\[1mm]
Cada error de hoy, bien observado, se convierte en pieza del camino que pisarás mañana.
\end{softbox}

\small
\begin{tabularx}{\textwidth}{>{\bfseries}p{3.6cm}Y}
\rowcolor{milcGradeDark!12}Nombre completo: & \linead{10cm}\\
Fecha: & \linead{4cm} \quad Curso/grupo: \linead{4cm}\\
Mi firma: & \linead{9cm}\\
\end{tabularx}
\normalsize

\begin{infoband}{milcGradeDark}
\textbf{Para la próxima sesión.} Llega con tu análisis auditado y la reflexión sobre límites. En la sesión 7 vas a aprender a escribir \textbf{informe comercial para audiencia real}: propuesta concreta dirigida a alguien que pueda decidir. El análisis de hoy es la materia; el informe comercial de mañana es la propuesta basada en datos.
\end{infoband}

\end{document}
